%%%%%%%%%%%%%%%%%%%%%%%%%%%%%%%%%%%%%%%%%%%%%%%%%%%%%%%%%%%%%%%%%%%%%%%%%%%%%%%
% CHAPTER 67: PHYSICAL ORIGINS
%%%%%%%%%%%%%%%%%%%%%%%%%%%%%%%%%%%%%%%%%%%%%%%%%%%%%%%%%%%%%%%%%%%%%%%%%%%%%%%

\chapter{Physical origins}
\label{chap:physical-origins}
\label{rem:programme-origins}


String amplitudes, gauge theory anomalies, and conformal blocks are translated into the bar-cobar formalism of Part~I; physical identifications are conjectural unless theorem-labeled.
\section{4d/2d correspondence algebras}
\label{sec:4d2d}

\subsection{The 4d/2d correspondence (BLLPRR)}

\begin{construction}[4d/2d correspondence (BLLPRR)]\label{constr:4d2d}
Starting from a 4d $\mathcal{N} = 2$ superconformal field theory $\mathcal{T}_{4d}$:
\begin{enumerate}[label=(\roman*)]
\item Select the \emph{Schur supercharge} $\mathbb{Q} = Q_1^- + \widetilde{S}^{1\dot{-}}$, a specific nilpotent supercharge in the $\mathcal{N}=2$ superconformal algebra ($\mathbb{Q}^2 = 0$).
\item Pass to the cohomology $H^*(\mathbb{Q})$ of local operators inserted at a fixed plane $\mathbb{R}^2 \subset \mathbb{R}^4$.
\item The resulting algebra $\mathcal{A}[\mathcal{T}_{4d}] = H^*(\mathbb{Q}, \text{local ops})$ carries the structure of a 2d vertex algebra (chiral algebra).
\end{enumerate}

This construction, due to Beem--Lemos--Liendo--Peelaers--Rastelli--van~Rees~\cite{beem-et-al}, produces vertex algebras from 4d SCFTs.
It is distinct from the $\Omega$-background/Nekrasov partition-function
approach used in the AGT correspondence below.
\end{construction}

\begin{theorem}[4d/2d is \texorpdfstring{$\Einf$}{E-infinity}-chiral; \ClaimStatusProvedElsewhere{} \cite{beem-et-al}]\label{thm:4d2d-einf}
The 2d chiral algebra $\mathcal{A}[\mathcal{T}_{4d}]$ obtained from the 4d/2d
correspondence is an $\Einf$-chiral algebra (vertex algebra).

The OPE of the resulting chiral algebra is generically \emph{non-commutative}: for instance, the 4d/2d image of $SU(2)$ $\mathcal{N}=4$ SYM is the small $\mathcal{N}=4$ superconformal vertex algebra at $c = -9$, which contains $\widehat{\mathfrak{sl}}_2$ at level $k = -3/2$ (an admissible, non-integral level; determined by $c = 6k$) as its R-symmetry current subalgebra.  The $\Einf$ structure arises from the full vertex algebra axioms (locality, associativity of OPE), not from commutativity.
\end{theorem}

\begin{example}[Class \texorpdfstring{$\mathcal{S}$}{S}]\label{ex:class-s}
For class $\mathcal{S}$ theories of type $A_{n-1}$ on a Riemann surface $C$:
\[
\mathcal{A}[\mathcal{T}_{A_{n-1}}(C)] = \mathcal{W}^{-n}(\mathfrak{sl}_n)
\]
the W-algebra associated to $\mathfrak{sl}_n$ at the critical level $k = -n = -h^\vee$. At this special level the Sugawara construction degenerates and $\mathcal{W}^{-h^\vee}(\mathfrak{sl}_n)$ coincides with the Feigin--Frenkel center $\mathfrak{z}(\widehat{\mathfrak{g}})$, a commutative vertex algebra (see Feigin--Frenkel \cite{FF} for the center isomorphism and Arakawa \cite{Ara15} for related rationality results). The dependence on the Riemann surface $C$ and its pants decomposition appears at the level of conformal blocks and modules, not the algebra itself.
\end{example}


\section{Non-commutative Chern--Simons theory}
\label{sec:nc-chern-simons}
\noindent\emph{All claims in this section are \ClaimStatusConjectured.}

\subsection{Chern--Simons as source of chiral algebras}

\begin{construction}[CS boundary algebra]\label{constr:cs-boundary}
3d Chern--Simons theory with gauge group $G$ at level $k$ on $M^3 = \Sigma \times \bR_+$ 
produces:
\begin{enumerate}[label=(\roman*)]
\item On the boundary $\Sigma$: WZW model at level $k$.
\item The chiral algebra is the affine Kac--Moody algebra $\hat{\mathfrak{g}}_k$.
\end{enumerate}
\end{construction}

\begin{conjecture}[Non-commutative CS; \ClaimStatusConjectured]\label{conj:nc-cs}
Noncommutative Chern--Simons theory on $\bR^3_\theta$ produces
$\Eone$-chiral algebras:
\begin{enumerate}[label=(\roman*)]
\item The boundary theory is a non-commutative WZW model.
\item The OPE is R-twisted with $R$ depending on the non-commutativity parameter $\theta$.
\item At $\theta = 0$, we recover the standard $\Einf$-chiral Kac--Moody.
\end{enumerate}
\end{conjecture}

\begin{remark}[Scope]
The $R$-twisted OPE structure is consistent with Definition~\ref{def:chiral-ass-operad}, but deriving it from non-commutative CS requires Seiberg--Witten deformation quantization~\cite{SW99}, which lies outside this monograph.  The remaining targets belong to MC4/MC5 (Chapter~\ref{chap:concordance}).  Contributing to Conjecture~\ref{conj:master-bv-brst}.
\end{remark}


\subsection{\texorpdfstring{Noncommutative geometry and associative / topological $\mathsf{E}_1$ structures}{Noncommutative geometry and associative / topological E1 structures}}

\begin{construction}[NC torus as associative deformation]
\label{constr:nc-torus-e1}
\index{noncommutative torus!topological $E_1$ structure@topological $\mathsf{E}_1$ structure}
The noncommutative torus $T^2_\theta$ ($\theta \in \mathbb{R} \setminus \mathbb{Q}$)
carries a natural associative deformation, equivalently a topological
$\Eone$-algebra structure.  The smooth algebra
$C^\infty(T^2_\theta)$ is generated by unitaries $U, V$ subject to $VU = e^{2\pi i \theta} UV$,
and the Moyal star product
\[
(f \star_\theta g)(x,y) = \sum_{m,n} \hat{f}(m,n)\,\hat{g}(m',n')\,
e^{i\pi\theta(mn'-m'n)}\,e^{2\pi i((m+m')x + (n+n')y)}
\]
is associative but noncommutative, yielding a topological
$\Eone$-structure.  What is \emph{not} constructed here is a
Beilinson--Drinfeld $\Eone$-chiral algebra on a curve.  The object at
hand is the associative deformation of the Poisson structure
$\{x,y\} = \theta$ on $T^2$.
\end{construction}

\begin{conjecture}[Koszul duality as Morita equivalence; \ClaimStatusConjectured]
\label{conj:koszul-morita}
\index{Morita equivalence!Koszul duality interpretation}
$\Eone$-Koszul duality for noncommutative tori exchanges the deformation parameter:
\[
\mathcal{A}_\theta^! \;\cong\; \mathcal{A}_{1/\theta}.
\]
This coincides with Rieffel's Morita equivalence $T^2_\theta \sim_{\mathrm{Morita}} T^2_{1/\theta}$
for irrational $\theta$, reinterpreted as $\Eone$-Koszul duality.
The bar complex $\barB(\mathcal{A}_\theta)$ computes the cyclic homology
$\mathrm{HC}_*(\mathcal{A}_\theta)$, which by Connes' theorem equals the de~Rham
cohomology $H^*_{\mathrm{dR}}(T^2)$ --- independent of $\theta$ (topological invariance
of cyclic homology for smooth deformations).
\end{conjecture}

\begin{remark}[Scope]
Rieffel Morita equivalence and Connes' $\mathrm{HC}_*(A_\theta)$ computation are established; the missing step is constructing the topological $\Eone$ bar complex for $\mathcal{A}_\theta$ and verifying its cohomology recovers the Koszul dual at $1/\theta$, which lies outside the algebraic-curve setting.  Contributing to Conjecture~\ref{conj:master-bv-brst}.
\end{remark}


\section{Gauge theory and D-branes}
\label{sec:gauge-d-branes}
\subsection{D-brane vertex algebras}

\begin{construction}[Open string VOA]\label{constr:open-string-voa}
Consider open strings ending on D-branes in type IIB string theory:
\begin{enumerate}[label=(\roman*)]
\item The worldsheet is a disk $D$ with boundary conditions.
\item The boundary CFT gives rise to a chiral algebra.
\item For D-branes in a Calabi--Yau, the algebra depends on the D-brane configuration.
\end{enumerate}
\end{construction}

\begin{conjecture}[D-brane algebras are \texorpdfstring{$\Eone$}{E1}; \ClaimStatusConjectured]\label{conj:dbrane-e1}
Open string vertex algebras on non-trivial D-brane configurations are generically 
$\Eone$-chiral algebras:
\begin{enumerate}[label=(\roman*)]
\item The Chan--Paton factors introduce non-commutativity.
\item The OPE has a matrix structure: $\Phi^{ij}(z) \Psi^{k\ell}(w) \sim \delta^{jk} \cdots$.
\item Commutativity of the OPE fails due to the matrix ordering (the OPE is no longer symmetric under exchange of operators, yielding an $\Eone$- rather than $\Einf$-structure).
\end{enumerate}
\end{conjecture}

\begin{remark}[Scope]
Matrix-valued OPE from Chan--Paton factors is compatible with $\Eone$-chiral structure (Definition~\ref{def:chiral-ass-operad}), but deriving it from open string field theory~\cite{Zwi93} lies outside this monograph.  The remaining targets follow the MC4/MC5 dependency order (Chapter~\ref{chap:concordance}).  Contributing to Conjecture~\ref{conj:master-bv-brst}.
\end{remark}


\subsection{\texorpdfstring{D-brane categories from $\Eone$-module Koszul duality}{D-brane categories from E1-module Koszul duality}}

\begin{conjecture}[HMS from \texorpdfstring{$\Eone$}{E1}-module Koszul duality; \ClaimStatusConjectured]
\label{conj:hms-koszul}
\index{homological mirror symmetry!Koszul duality interpretation}
\index{D-brane!module Koszul duality}
In string theory, D-branes on a target space $X$ form the derived category
$D^b(\mathrm{Coh}(X))$ (B-branes) or the Fukaya category $\mathrm{Fuk}(X)$
(A-branes).  For a chiral algebra $\mathcal{A}$ on a curve $\Sigma$, the
module category $\mathrm{Mod}(\mathcal{A})$ is the mathematical incarnation
of the category of branes.

Under $\Eone$-Koszul duality, the module bar-cobar functor
$\Phi\colon \mathrm{Mod}(\mathcal{A}) \xrightarrow{\sim}
\mathrm{CoMod}(\mathcal{A}^!)$
of Theorem~\ref{thm:e1-module-koszul-duality} exchanges module and comodule
categories.  When $X$ is a Calabi--Yau surface and $\mathcal{A}$ is the chiral
algebra of the sigma model on $X$, this functor should recover homological
mirror symmetry (Kontsevich):
\[
D^b\bigl(\mathrm{Coh}(X)\bigr)
\;\simeq\;
D^b\bigl(\mathrm{Fuk}(X^\vee)\bigr)
\]
for a mirror pair $(X, X^\vee)$.  In this interpretation, the bar-cobar
adjunction is conjecturally interpreted as the passage from B-branes to A-branes: the bar construction
$\barB(\mathcal{A})$ computes the Ext-algebra of the diagonal brane,
and the cobar construction recovers the endomorphism algebra of the dual
collection on $X^\vee$.
\end{conjecture}

\begin{remark}[Scope]
Theorem~\ref{thm:e1-module-koszul-duality} stands independently; the HMS interpretation requires identifying chiral algebra modules with brane categories (boundary CFT input) and passing from the abstract functor $\Phi$ to derived category equivalences (tilting generation).  Contributing to Conjecture~\ref{conj:master-bv-brst}.
\end{remark}


\section{AGT correspondence connections}
\label{sec:agt}

\begin{theorem}[AGT correspondence; \ClaimStatusProvedElsewhere{} \cite{AGT09}]\label{thm:agt}
(Alday--Gaiotto--Tachikawa) For 4d $\mathcal{N} = 2$ $SU(2)$ gauge theory with 
$N_f = 4$ fundamental hypermultiplets:
\[
Z_{\mathrm{Nekrasov}}(\epsilon_1, \epsilon_2, a; q) = 
\langle V_{\alpha_1}(z_1) \cdots V_{\alpha_4}(z_4) \rangle_{\mathrm{Liouville}}
\]
where:
\begin{itemize}
\item LHS: Nekrasov partition function
\item RHS: 4-point conformal block in Liouville CFT
\item $c = 1 + 6Q^2$, $Q = b + b^{-1}$, $b^2 = \epsilon_1/\epsilon_2$ (dimensionless ratio of $\Omega$-background parameters)
\end{itemize}
\end{theorem}

\begin{corollary}[Virasoro from gauge theory; \ClaimStatusProvedElsewhere{} \cite{AGT09}]\label{cor:virasoro-gauge}
The Virasoro algebra (an $\Einf$-chiral algebra) arises from the gauge theory 
computation. The central charge $c$ is determined by the $\Omega$-background parameters.
\end{corollary}


\begin{conjecture}[\texorpdfstring{$q$}{q}-AGT; \ClaimStatusConjectured]\label{conj:q-agt}
The 5d lift of AGT (adding a circle) gives:
\[
Z_{\mathrm{5d Nekrasov}}(q, t) = \text{(conformal blocks of } \mathcal{W}_{q,t})
\]
where $\mathcal{W}_{q,t}$ is the quantum W-algebra, an $\Eone$-chiral algebra.
\end{conjecture}

\begin{remark}[Scope]
$\mathcal{W}_{q,t}$ as an $\Eone$-chiral algebra is well-defined (Frenkel--Reshetikhin); the equality with 5d Nekrasov partition functions~\cite{AY10} is a physics conjecture beyond the methods of this monograph.  Contributing to Conjecture~\ref{conj:master-bv-brst}.
\end{remark}

